% research-writing preamble (research-writing skill; see its SKILL.md). Edit the document, not this file:
% project-wide changes go in a local preamble-local.tex, which main.tex inputs after this one.
\usepackage{iftex}
\usepackage{amsmath}
\ifLuaTeX\usepackage{fontspec}\usepackage{libertinus}\else\usepackage[T1]{fontenc}\usepackage{lmodern}\usepackage{amssymb}\fi
\usepackage[a4paper,margin=25mm]{geometry}
\usepackage{microtype}
\usepackage{graphicx}
\usepackage{booktabs,tabularx}
\usepackage[detect-all]{siunitx}
\usepackage{xcolor}
\usepackage[font=small,labelfont=bf,skip=4pt]{caption}
\usepackage{subcaption}
\usepackage[section]{placeins} % a figure never floats past the end of its section
\usepackage[backend=biber,style=authoryear,maxcitenames=2,uniquename=false,uniquelist=false]{biblatex}
\usepackage[hidelinks]{hyperref}
\usepackage[nameinlink,capitalise,noabbrev]{cleveref}
\usepackage{enumitem}
\setlist{nosep,leftmargin=*}
\setlength{\parskip}{0.4\baselineskip}\setlength{\parindent}{0pt}
% floats: let figures sit on text pages and near their first citation
\renewcommand{\topfraction}{0.85}\renewcommand{\textfraction}{0.12}\renewcommand{\floatpagefraction}{0.75}
\setcounter{topnumber}{2}\setcounter{totalnumber}{3}

% ---- the skill's marks (checked by wr.py check)
% \claim{C3}: this paragraph carries claim C3 of ledger.md (prints nothing)
\newcommand{\claim}[1]{}
% \fact{F3}: the next sentence states ledger fact F3 (prints nothing; wr.py checks each fact is stated once)
\newcommand{\fact}[1]{}
% \val{name}: a named value from values/*.json (values.tex is generated by wr.py values); robust, so it also works in a
% \caption and in the In short box. An undefined name prints a bold ??name and a package warning. Names go only into
% \csname, so underscores are safe; the undefined branch prints the name with \detokenize.
\makeatletter
\newcommand{\wrsetval}[2]{\expandafter\def\csname wrval@#1\endcsname{#2}}
\DeclareRobustCommand{\val}[1]{%
  \ifcsname wrval@#1\endcsname\csname wrval@#1\endcsname
  \else\textbf{??\detokenize{#1}}\PackageWarning{wr}{value \detokenize{#1} undefined}\fi}
% in PDF bookmarks (a \val in a section title) the value as plain text
\newcommand{\wrvalof}[1]{\ifcsname wrval@#1\endcsname\csname wrval@#1\endcsname\else ??\detokenize{#1}\fi}
\pdfstringdefDisableCommands{\let\val\wrvalof}
\makeatother
% \valw{name}, \Valw{name}: the value as a number word ("six", "Six" at a sentence start) for whole numbers 0-20, so
% the words around a result follow it too; any other value prints as \val does
\makeatletter
\newcommand{\wrsetvalw}[3]{\expandafter\def\csname wrvalw@#1\endcsname{#2}\expandafter\def\csname wrvalW@#1\endcsname{#3}}
\DeclareRobustCommand{\valw}[1]{\ifcsname wrvalw@#1\endcsname\csname wrvalw@#1\endcsname\else\val{#1}\fi}
\DeclareRobustCommand{\Valw}[1]{\ifcsname wrvalW@#1\endcsname\csname wrvalW@#1\endcsname\else\val{#1}\fi}
\pdfstringdefDisableCommands{\let\valw\wrvalof\let\Valw\wrvalof}
\makeatother
% \valw[none]{name}: the word in brackets when the value is 0 ("none of the 13 pass"), else as \valw;
% \valp{name}{draw}{draws}: the singular when the value is 1, else the plural, so "one draw is" follows the value
\makeatletter
\newcommand{\wrsetvalc}[2]{\expandafter\def\csname wrvalc@#1\endcsname{#2}}
\def\wr@z{z}\def\wr@o{o}
\newcommand{\wr@cls}[1]{\edef\wr@c{\ifcsname wrvalc@#1\endcsname\csname wrvalc@#1\endcsname\else x\fi}}
\newcommand{\wr@word}[2]{\ifcsname #1@#2\endcsname\csname #1@#2\endcsname\else\val{#2}\fi}
\DeclareRobustCommand\valw[2][]{\wr@cls{#2}\def\wr@a{#1}%
  \ifx\wr@c\wr@z\ifx\wr@a\@empty\wr@word{wrvalw}{#2}\else#1\fi\else\wr@word{wrvalw}{#2}\fi}
\DeclareRobustCommand\Valw[2][]{\wr@cls{#2}\def\wr@a{#1}%
  \ifx\wr@c\wr@z\ifx\wr@a\@empty\wr@word{wrvalW}{#2}\else#1\fi\else\wr@word{wrvalW}{#2}\fi}
\DeclareRobustCommand\valp[3]{\wr@cls{#1}\ifx\wr@c\wr@o#2\else#3\fi}
\def\wr@pdfw#1#{\wrvalof}\def\wr@pdfp#1#2#3{#3}
\pdfstringdefDisableCommands{\let\valw\wr@pdfw\let\Valw\wr@pdfw\let\valp\wr@pdfp}
\makeatother
% \figref{fig:x}: cite a figure; records the page of the first citation for wr.py (page of figure vs of its first mention).
% Robust, so it also works inside a \caption (a moving argument; a plain \newcommand is fatal there).
\makeatletter
\DeclareRobustCommand{\figref}[1]{\cref{#1}\protected@write\@auxout{}{\string\wrfigcite{#1}{\thepage}}}
\DeclareRobustCommand{\Figref}[1]{\Cref{#1}\protected@write\@auxout{}{\string\wrfigcite{#1}{\thepage}}}
\newcommand{\wrfigcite}[2]{}
\makeatother
% \status{established|plausible|open}: the evidence label of a claim, printed as a small tag
\newcommand{\status}[1]{\textsuperscript{\textsf{\scriptsize[#1]}}}
% a boxed summary for the first page of a write-up
\usepackage[most]{tcolorbox}
\newtcolorbox{inshort}{colback=black!4,colframe=black!4,boxrule=0pt,arc=0pt,left=6pt,right=6pt,top=4pt,bottom=4pt,
  before upper={\textbf{In short.}\enskip}}

% \lead{The point.} a bold run-in head that opens a paragraph and names its point
\newcommand{\lead}[1]{\textbf{#1}\enspace}

% \supplementary: after \printbibliography; a new page titled Supplementary Information, sections, figures and
% tables numbered S1, S2 (wr.py check exempts figures after it from the near-first-citation rule)
\makeatletter
\newcommand{\supplementary}{%
  \clearpage
  \begin{center}{\LARGE\bfseries Supplementary Information}\end{center}\vspace{1em}
  \setcounter{section}{0}\setcounter{figure}{0}\setcounter{table}{0}
  \renewcommand{\thesection}{S\arabic{section}}\renewcommand{\thefigure}{S\arabic{figure}}%
  \renewcommand{\thetable}{S\arabic{table}}%
  \renewcommand{\theHsection}{S\arabic{section}}\renewcommand{\theHfigure}{S\arabic{figure}}%
  \renewcommand{\theHtable}{S\arabic{table}}%
}
\makeatother
\IfFileExists{values.tex}{\input{values}}{}
